\section{研究倫理}
\label{sec:ethics}

本研究は，公開されている脆弱性情報のみを対象とし，脆弱性の新規発見や実システムへの攻撃を伴わない．
調査に用いたデータは，NVDが公開するJSON 2.0データフィード，VulnCheck NVD++のAPI，および各ベンダ・プロジェクトが公開するセキュリティアドバイザリである．
いずれも一般に公開された情報であり，個人情報や非公開情報は含まない．

VulnCheck NVD++のデータ利用にあたっては，提供元が定める利用条件に従い，当該データがVulnCheck NVD++に由来することを明示する．
また各ベンダの公開情報の取得にあたっては，提供元サーバに過度な負荷をかけないよう取得間隔を制御する．

6章で提案するLLMの利用においては，LLMに入力するのは公開されたアドバイザリの本文のみであり，非公開情報を入力しない．
また，LLMの出力を検証せずに脆弱性データベースへ投入しない設計としており，誤った脆弱性情報が利用者に届くリスクを低減している．

本稿で報告する集計値は，公開データに対する機械的な集計であり，特定のベンダや組織の評価・非難を意図するものではない．
NVDの \texttt{Deferred} は運用方針の表明として公開されているものであり，本稿はそれを所与の条件として利用者側の対応を論じている．
